\Needspace{8\baselineskip}
\Article{194}
Da die Formen, die wir oben assoziiert genannt haben (Artikel 187, 6), stets uneigentlich äquivalent sind (Artikel 159), ist klar: Sind die reduzierten Formen $F$, $f$ uneigentlich äquivalent und ist die zu $F$ assoziierte Form $G$, so sind $f$, $G$ eigentlich äquivalent, und daher ist $G$ in der Periode von $f$ enthalten. Sind also die Formen $F$, $f$ sowohl eigentlich als auch uneigentlich äquivalent, so müssen sowohl $F$ als auch $G$ in der Periode von $f$ vorkommen. Diese Periode ist deshalb zu sich selbst assoziiert und enthält zwei ambige Formen (Artikel 187, 7). Das bestätigt auf schöne Weise den Satz in Artikel 165, aufgrund dessen wir bereits sicher sein konnten, dass es eine ambige Form gibt, die zu $F$, $f$ äquivalent ist.

\Needspace{8\baselineskip}
\Article{195}
\textsc{Aufgabe.} \textit{Für zwei beliebige gegebene Formen $\Phi$, $\varphi$ mit derselben Determinante entscheiden, ob sie äquivalent sind oder nicht.}

\textit{Lösung.} Wir suchen zwei reduzierte Formen $F$, $f$, die zu den gegebenen Formen $\Phi$, $\varphi$ der Reihe nach eigentlich äquivalent sind (Artikel 183). Je nachdem, ob diese nur eigentlich, nur uneigentlich, auf beide Arten oder auf keine Art äquivalent sind, sind auch die gegebenen Formen nur eigentlich, nur uneigentlich, auf beide Arten oder gar nicht äquivalent. Wir bestimmen die Periode einer der beiden reduzierten Formen, \textit{z.B.} die von $f$. Kommt die Form $F$ in dieser Periode vor, aber nicht zugleich die zu $F$ assoziierte Form, so liegt offenbar der \textit{erste} Fall vor; kommt dagegen diese assoziierte Form vor, aber nicht $F$ selbst, der \textit{zweite}; kommen beide vor, der \textit{dritte}; kommt keine vor, der \textit{vierte}.

\textit{Beispiel.} Gegeben seien die Formen $(129,92,65)$, $(42,59,81)$ mit der Determinante 79. Als dazu eigentlich äquivalente reduzierte Formen finden wir $(10,7,-3)$, $(5,8,-3)$. Die Periode der ersten ist $(10,7,-3)$, $(-3,8,5)$, $(5,7,-6)$, $(-6,5,9)$, $(9,4.\,-7)$, $(-7,3,10)$. Da darin die Form $(5,8,-3)$ selbst nicht vorkommt,
\SourceLeaf{190}
wohl aber ihre assoziierte Form $(-3,8,5)$, schließen wir, dass die gegebenen Formen nur uneigentlich äquivalent sind.

Teilen wir alle reduzierten Formen mit einer gegebenen Determinante wie oben (Artikel 187, 5) in Perioden $P$, $Q$, $R$ usw. ein und wählen wir aus jeder Periode eine beliebige Form, aus $P$ die Form $F$, aus $Q$ die Form $G$, aus $R$ die Form $H$ usw., so können unter diesen Formen $F,G,H$ usw. keine zwei eigentlich äquivalent sein. Jede andere Form mit derselben Determinante ist aber zu einer von ihnen eigentlich äquivalent, und zwar zu \textit{genau einer}. Offenbar lassen sich also \textit{alle Formen mit dieser Determinante in ebenso viele Klassen einteilen, wie es Perioden gibt}, indem wir die zu $F$ eigentlich äquivalenten Formen in die erste Klasse, die zu $G$ eigentlich äquivalenten in die zweite usw. aufnehmen. So sind alle Formen in derselben Klasse eigentlich äquivalent; Formen aus verschiedenen Klassen können dagegen nicht eigentlich äquivalent sein. Wir halten uns aber hier nicht bei diesem Thema auf, das wir unten ausführlicher behandeln werden.

\Needspace{8\baselineskip}
\Article{196}
\textsc{Aufgabe.} \textit{Für zwei gegebene eigentlich äquivalente Formen $\Phi$, $\varphi$ eine eigentliche Transformation der einen in die andere finden.}

\textit{Lösung.} Nach der Methode in Artikel 183 können wir zwei Folgen von Formen
\[
\Phi,\Phi',\Phi''\mathinner{\ldotp\ldotp\ldotp}\Phi^n\quad\text{und}\quad\varphi,\varphi',\varphi''\mathinner{\ldotp\ldotp\ldotp}\varphi^\nu
\]
finden, in denen jede Form zu ihrer Vorgängerin eigentlich äquivalent ist und deren letzte Formen $\Phi^n$, $\varphi^\nu$ reduziert sind. Da $\Phi$, $\varphi$ nach Voraussetzung eigentlich äquivalent sind, muss $\Phi^n$ in der Periode von $\varphi^\nu$ enthalten sein. Sei $\varphi^\nu=f$, und ihre Periode bis zur Form $\Phi^n$ sei
\[
f,\ f',\ f''\mathinner{\ldotp\ldotp\ldotp} f^{m-1},\ \Phi^n
\]
sodass $\Phi^n$ in dieser Periode den Index $m$ hat. Wir bezeichnen die Formen, die den assoziierten Formen von
\begin{center}
$\Phi,\ \Phi',\ \Phi''\mathinner{\ldotp\ldotp\ldotp}\Phi^n$\quad entgegengesetzt sind, mit $\Psi,\ \Psi',\ \Psi''\mathinner{\ldotp\ldotp\ldotp}\Psi^n$ der Reihe nach\HistoricalNote{*)}{Es entsteht also $\Psi$ aus $\Phi$, indem wir das erste und letzte Glied vertauschen und dem mittleren das entgegengesetzte Vorzeichen geben; ebenso bei den anderen.}.
\end{center}
Dann ist in der Folge
\SourceLeaf{191}
\[
\varphi,\ \varphi',\ \varphi''\mathinner{\ldotp\ldotp\ldotp} f,\ f',\ f''\mathinner{\ldotp\ldotp\ldotp} f^{m-1},\ \Psi^{n-1},\ \Psi^{n-2}\mathinner{\ldotp\ldotp\ldotp}\Psi,\ \Phi
\]
jede Form ihrer Vorgängerin nach rechts benachbart. Daher können wir nach Artikel 177 eine eigentliche Transformation der ersten Form $\varphi$ in die letzte $\Phi$ finden. Bei den übrigen Formen der Folge ist das leicht zu sehen; für $f^{m-1}$, $\Psi^{n-1}$ beweisen wir es wie folgt. Sei
\[
\begin{gathered}f^{m-1}=(g,h,i);\quad f^m\text{ oder }\Phi^n=(g',h',i');\\ \Phi^{n-1}=(g'',h'',i'')\end{gathered}
\]
Die Form $(g',h',i')$ ist sowohl der Form $(g,h,i)$ als auch der Form $(g'',h'',i'')$ nach rechts benachbart. Daher gilt $i=g'=i''$ und $-h\equiv h'\equiv-h''\ (\mathrm{mod.}\ i\text{ oder }g'\text{ oder }i'')$. Somit ist die Form $(i'',-h'',g'')$, \textit{d.h.} die Form $\Psi^{n-1}$, der Form $(g,h,i)$, \textit{d.h.} der Form $f^{m-1}$, nach rechts benachbart.

Sind die Formen $\Phi$, $\varphi$ uneigentlich äquivalent, so ist $\varphi$ zu der Form eigentlich äquivalent, die $\Phi$ entgegengesetzt ist. Wir können also eine eigentliche Transformation von $\varphi$ in die zu $\Phi$ entgegengesetzte Form finden. Erfolgt diese durch die Substitution $\alpha,\beta,\gamma,\delta$, so sehen wir leicht, dass $\varphi$ uneigentlich in $\Phi$ selbst transformiert wird, und zwar durch die Substitution $\alpha,-\beta,\gamma,-\delta$.

Damit ist auch klar, dass wir für die Formen $\Phi$, $\varphi$, wenn sie sowohl eigentlich als auch uneigentlich äquivalent sind, \textit{zwei} Transformationen finden können, eine eigentliche und eine uneigentliche.

\Needspace{12\baselineskip}
\textit{Beispiel.} Gesucht ist eine uneigentliche Transformation der Form $(129,92,65)$ in $(42,59,81)$, die wir im vorhergehenden Artikel als dazu uneigentlich äquivalent erkannt haben. Wir müssen also zuerst eine eigentliche Transformation von $(129,92,65)$ in $(42,-59,81)$ suchen. Dazu bestimmen wir die folgende Formenfolge: $(129,92,65)$, $(65,-27,10)$, $(10,7,-3)$, $(-3,8,5)$, $(5,22,81)$, $(81,59,42)$, $(42,-59,81)$. Daraus leiten wir die eigentliche Transformation $-47,56,73,-87$ her, durch die $(129,92,65)$ in $(42,-59,81)$ übergeht; durch die uneigentliche Transformation $-47,\allowbreak-56,\allowbreak73,\allowbreak87$ geht sie also in $(42,59,81)$ über.
